\documentclass[11pt]{article}
\usepackage[margin=1in]{geometry}
\usepackage{amsmath}
\usepackage{xurl}
\usepackage{enumitem}
\usepackage{titlesec}
\usepackage[colorlinks=true,linkcolor=blue,citecolor=blue,urlcolor=blue]{hyperref}

\titleformat{\subsection}{\normalfont\bfseries\large}{\thesubsection}{1em}{}
\setlist[itemize]{leftmargin=1.4em, itemsep=0.2em, topsep=0.2em}

\title{Literature Review: Cannabis Entry Restrictions and the Cross-Border Demand Premium\\
\large Supporting document -- paper-by-paper summary and contribution mapping}
\author{}
\date{Compiled \today}

\begin{document}
\maketitle

\section*{Purpose and scope}
This document lists every paper identified in a comprehensive literature search on (1) cannabis retail entry
restrictions (license caps, quotas, local bans, moratoria) and (2) cross-border cannabis demand, and states
plainly what our own paper contributes relative to each one. Full search methodology and coverage caveats are
in Section~\ref{sec:methodology}. Two entries (marked \textbf{[UNVERIFIED]}) were identified by title only and
have not been confirmed against the actual paper text -- do not cite them without independent verification.

\section{Papers studying cannabis entry restrictions}

\subsection{Thomas, D. (2019) -- ``License Quotas and the Inefficient Regulation of Sin Goods: Evidence from
the Washington Recreational Marijuana Market''}
\begin{itemize}
\item \textbf{Venue:} SSRN working paper 3312960 (no journal publication identified as of this search)
\item \textbf{Sample and years:} Washington State, retail licensing 2014--2017
\item \textbf{Identification/method:} Structural model of Washington's capped, weighted-lottery license
allocation system
\item \textbf{Main findings:} Free entry would raise total surplus by roughly 18 percent relative to
Washington's actual capped and partly randomized allocation; decomposes the loss into geographic
misallocation (58 percent of the total) and the randomness of the lottery itself (13 percent)
\item \textbf{Border-demand findings:} None -- no cross-border or neighboring-state demand variable in the
design
\item \textbf{Our contribution relative to this paper:} Thomas identifies a real welfare cost of Washington's
cap but has no demand shock to test it against -- the model is single-state and cannot speak to how an
external, plausibly exogenous demand shock (a neighboring state legalizing) would be absorbed differently
under a cap versus free entry. We supply exactly that: a multi-state design with entry-regime variation and a
time-varying demand shock, letting us test whether the type of distortion Thomas documents structurally also
shows up, in reduced form, specifically on the quantity margin when a demand shock hits. Our \$490 million
illustrative revenue-gap estimate is a different (larger, cruder) quantity than Thomas's 18 percent surplus
loss -- we discuss this comparison explicitly in our own paper -- but the two results corroborate each other's
qualitative conclusion using entirely different data and methods.
\end{itemize}

\subsection{Hollenbeck, B. and Uetake, K. (2021) -- ``Taxation and Market Power in the Legal Marijuana
Industry''}
\begin{itemize}
\item \textbf{Venue:} \emph{RAND Journal of Economics}, Vol.\ 52(3), pp.\ 559--595
\item \textbf{Sample and years:} Washington State, near-universe of administrative transaction data
\item \textbf{Identification/method:} Structural IO model of demand and markups using firm-level transaction
records
\item \textbf{Main findings:} Washington's strict entry cap produces retailers with substantial market power;
the state's 37 percent excise tax is borne mostly by consumers rather than firms, despite that market power
\item \textbf{Border-demand findings:} None -- single state, no cross-border variation
\item \textbf{Our contribution relative to this paper:} Hollenbeck and Uetake show capped entry supports
markups using rich transaction-level data for one state. Our reduced-form, five-state result is consistent
with and extends their finding: we show that under capped entry, revenue per store stays essentially flat
regardless of border exposure (a precise null, $p=0.959$) while the quantity margin (dispensary density,
revenue per capita) is sharply suppressed -- exactly the pattern their markup/rent story would predict if
incumbents extract the demand shock as price rather than losing it to new entrants. We provide corroborating
reduced-form evidence for their mechanism from an entirely different empirical setting (a demand shock,
not their tax-incidence design), and extend it beyond Washington to four additional states.
\end{itemize}

\subsection{Payan, D., Brown, P., and Song, A. (2021) -- county-level census of Colorado and Washington local
marijuana ordinances}
\begin{itemize}
\item \textbf{Venue:} Already cited and verified in our own manuscripts
\item \textbf{Sample and years:} Colorado and Washington, county/local-jurisdiction ordinances as of April 2019
\item \textbf{Identification/method:} Descriptive census of local government marijuana retail ordinances
\item \textbf{Main findings:} 48.4 percent of Colorado counties ban marijuana retail facilities outright; only
23.4 percent allow every facility type without restriction
\item \textbf{Border-demand findings:} None -- purely institutional documentation, no market-outcome
estimation at all
\item \textbf{Our contribution relative to this paper:} We use Payan, Brown, and Song's documentation directly
as the evidentiary basis for coding Colorado as a capped-entry state despite its lack of a state-level
numerical cap (our manuscript's Institutional Background section credits them explicitly). Their paper stops
at documenting which counties restrict entry; ours goes on to estimate what that restriction actually does to
market outcomes, and specifically whether it changes how a cross-border demand shock is absorbed.
\end{itemize}

\subsection{Yates, D. and Speer, J. (2018) -- ``Over and Under-Regulation in the Colorado Cannabis Industry:
A Data-Analytic Perspective''}
\begin{itemize}
\item \textbf{Venue:} \emph{International Journal of Drug Policy}, Vol.\ 59, pp.\ 63--66
\item \textbf{Sample and years:} Colorado, qualitative/descriptive
\item \textbf{Identification/method:} Policy-analytic commentary, not a causal empirical design
\item \textbf{Main findings:} Argues inconsistent over- and under-regulation, including retail licensing
policy, has constrained Colorado's legal industry and left room for illegal-market activity
\item \textbf{Border-demand findings:} None
\item \textbf{Our contribution relative to this paper:} Yates and Speer raise the concern qualitatively that
Colorado's regulatory patchwork distorts market outcomes; we provide the first quantitative, causally
identified test of that concern, using variation in local entry restriction and a specific, measurable
mechanism (cross-border demand) rather than a general regulatory-quality argument.
\end{itemize}

\subsection{Escudero, M. (\textasciitilde 2018) -- ``Rule-of-Thumb Pricing: Retail Cannabis in Washington
State'' \textbf{[UNVERIFIED -- confirmed by title/abstract only]}}
\begin{itemize}
\item \textbf{Venue:} Northwestern University job-market paper
\item \textbf{Sample and years:} Washington State (exact years not confirmed in this search)
\item \textbf{Identification/method:} Pricing-behavior model within Washington's capped-entry retail
structure (not independently confirmed beyond the title/abstract)
\item \textbf{Main findings:} Studies retailer pricing rules under capped market structure; specifics not
verified in this search
\item \textbf{Border-demand findings:} None apparent from title/abstract
\item \textbf{Our contribution relative to this paper:} If confirmed, this paper studies pricing behavior
within a single, capped regime with no cross-regime or cross-border comparison. Our design compares capped
against free-entry regimes directly and introduces the border-demand dimension entirely absent here. Verify
before citing.
\end{itemize}

\subsection{Kim, D., O'Connor, S., and Norwood, B. (\textasciitilde 2020) -- title suggests marijuana retail
deregulation and housing prices \textbf{[UNVERIFIED -- confirmed by title only, content not read]}}
\begin{itemize}
\item \textbf{Venue:} SSRN working paper 3537860
\item \textbf{Sample and years:} Not confirmed (likely Washington or Colorado)
\item \textbf{Identification/method:} Not confirmed; title implies a hedonic housing-price approach
\item \textbf{Main findings:} Not confirmed in this search
\item \textbf{Border-demand findings:} Not confirmed
\item \textbf{Our contribution relative to this paper:} Cannot be stated with confidence until the paper is
verified. If it studies entry deregulation's effect on housing prices, the outcome variable (housing values)
and likely absence of a cross-border demand mechanism would make it a different question from ours regardless
of what it finds. \textbf{Do not cite until independently verified.}
\end{itemize}

\subsection{Childs, C. et al.\ (2023) -- Saskatchewan private cannabis retail competition study}
\begin{itemize}
\item \textbf{Venue:} \emph{Canadian Public Administration}
\item \textbf{Sample and years:} Saskatchewan, Canada (not a U.S.\ state)
\item \textbf{Identification/method:} Competition-dimension analysis (location, pricing, customer experience)
of licensed private retailers
\item \textbf{Main findings:} Finds multidimensional competition among Saskatchewan's licensed urban cannabis
retailers
\item \textbf{Border-demand findings:} None
\item \textbf{Our contribution relative to this paper:} Different country, different institutional licensing
model (private retail under a national legal framework, not U.S.\ state-level prohibition/legalization
variation), no cross-border demand mechanism, and no entry-cap-versus-free-entry comparison. Our contribution
is essentially orthogonal to this paper's scope.
\end{itemize}

\section{Papers studying cross-border cannabis demand}

\subsection{Hansen, B., Miller, K., and Weber, C. (2020) -- ``Federalism, Partial Prohibition, and
Cross-Border Sales: Evidence from Recreational Marijuana''}
\begin{itemize}
\item \textbf{Venue:} \emph{Journal of Public Economics}, Vol.\ 187 (NBER Working Paper 23762, 2017)
\item \textbf{Sample and years:} Washington State retailers near the Oregon border, around Oregon's October
2015 retail launch
\item \textbf{Identification/method:} Difference-in-discontinuities design around Oregon's legalization date,
store-level monthly transaction data
\item \textbf{Main findings:} Washington retailers near the Oregon border lost roughly 36 percent of their
sales when Oregon's own retail market opened, direct evidence that the pre-existing premium those counties
enjoyed was a genuine captive cross-border demand effect
\item \textbf{Entry-restriction findings:} None -- single state (Washington, capped throughout), no
entry-regime variation tested
\item \textbf{Our contribution relative to this paper:} This is the closest existing paper to our own design
and the one our own methodology most directly builds on. Hansen, Miller, and Weber establish that cross-border
cannabis demand is real and economically large, but their setting has only one regulatory regime (Washington's
cap) and cannot ask whether a \emph{different} entry regime would have produced a different-sized effect. Our
central contribution is exactly that extension: a five-state design with genuine variation in entry
regulation, showing that the type of border effect they document is not regime-invariant -- it is
significantly larger under free entry than under capped entry for revenue per capita and dispensary density,
while remaining essentially fixed for revenue per store. Our per-capita magnitude (+196 percent) is larger
than but directionally consistent with the order of magnitude in their 36 percent sales-loss estimate; the
two are not directly comparable numbers (their estimate is a loss when a neighbor opens, ours is a premium
level), but both describe cross-border demand effects of several tens of percentage points or more, not a
marginal phenomenon.
\end{itemize}

\subsection{Hao, Z. and Cowan, B. (2020) -- ``The Cross-Border Spillover Effects of Recreational Marijuana
Legalization''}
\begin{itemize}
\item \textbf{Venue:} \emph{Economic Inquiry} (NBER Working Paper 23426)
\item \textbf{Sample and years:} Multi-state, recreational marijuana legalization (RML) states and their
neighbors
\item \textbf{Identification/method:} Difference-in-differences using RML legalization timing across states
\item \textbf{Main findings:} Recreational legalization in one state (e.g., Colorado, Washington) increases
marijuana possession arrests in neighboring nonlegal states, consistent with residents crossing the border to
purchase legally and consuming or being caught at home
\item \textbf{Entry-restriction findings:} None
\item \textbf{Our contribution relative to this paper:} Hao and Cowan study a public-safety/enforcement
outcome (arrests) in the \emph{nonlegal} neighboring state, not market outcomes in the legal state itself, and
do not test any entry-regulation moderator. Our outcomes (revenue, dispensary density) are the direct legal
cannabis market analogue of the demand shift they document indirectly through arrest data, and we add the
entry-regulation interaction entirely.
\end{itemize}

\subsection{Baggio, M., Chong, A., and Kwon, S. (2020) -- medical marijuana legalization and alcohol
substitution across state borders}
\begin{itemize}
\item \textbf{Venue:} Already cited and verified in our own manuscripts
\item \textbf{Sample and years:} U.S.\ states, medical marijuana legalization era
\item \textbf{Identification/method:} Border-county-pair design in the tradition of Dube, Lester, and Reich
(2010)
\item \textbf{Main findings:} Counties in states that legalized medical marijuana reduce alcohol purchases
relative to nearby counties just across the state line, evidence of cross-good substitution driven by legal
status differences at the border
\item \textbf{Entry-restriction findings:} None
\item \textbf{Our contribution relative to this paper:} Baggio, Chong, and Kwon study substitution to a
\emph{different} good (alcohol) as a consequence of cannabis legalization, using the same border-pair logic
our own design extends. We instead study cannabis's own market outcomes directly and add the entry-regulation
moderator, which their design does not include.
\end{itemize}

\subsection{Dube, A., Lester, T.\ W., and Reich, M.\ (2010) -- border-discontinuity minimum wage design}
\begin{itemize}
\item \textbf{Venue:} Already cited and verified in our own manuscripts
\item \textbf{Sample and years:} U.S.\ county pairs straddling state borders, minimum wage variation
\item \textbf{Identification/method:} Contiguous border-county-pair design (not cannabis)
\item \textbf{Main findings:} Establishes the border-discontinuity identification strategy as a way to net out
unobserved local economic conditions when comparing policy regimes across a state line
\item \textbf{Entry-restriction findings:} N/A -- not a cannabis paper
\item \textbf{Our contribution relative to this paper:} This is a methodological foundation, not a substantive
comparison point. We extend their static border-pair logic to a \emph{time-varying} border construction
(county exposure changes as neighboring states themselves legalize over our eleven-year panel) and apply it,
for the first time in this literature, jointly with an entry-regulation interaction.
\end{itemize}

\subsection{Friberg, R.\ et al.\ (2024) -- Norway COVID-19 border closures and cross-border shopping}
\begin{itemize}
\item \textbf{Venue:} Already cited and verified in our own manuscripts
\item \textbf{Sample and years:} Norway, COVID-19 border closure period
\item \textbf{Identification/method:} Natural experiment using COVID-19 border closures as an exogenous
shutoff of cross-border shopping
\item \textbf{Main findings:} Cross-border shopping measurably erodes domestic commodity tax revenue when the
border is open, evidence the underlying cross-jurisdiction demand-shopping mechanism is not specific to
cannabis or to the United States
\item \textbf{Entry-restriction findings:} None; not a licensing-cap setting at all
\item \textbf{Our contribution relative to this paper:} Friberg et al.\ study a blunt, binary policy
instrument (a full border closure) in a non-cannabis, non-U.S.\ context, with no entry-regulation margin.
We study a graduated regulatory instrument (entry caps of varying strictness) in the U.S.\ cannabis context,
and directly test whether that instrument moderates the underlying cross-border demand mechanism their paper
also documents in a different setting.
\end{itemize}

\subsection{Uruguay-to-Brazil cannabis cross-border spillover study (2025)}
\begin{itemize}
\item \textbf{Venue:} \emph{Journal of Quantitative Criminology}
\item \textbf{Sample and years:} Brazilian municipalities near the Uruguayan border, post-Uruguay
legalization
\item \textbf{Identification/method:} Difference-in-differences using distance from the Uruguay border
\item \textbf{Main findings:} Uruguay's national cannabis legalization affects drug seizures in nearby
Brazilian municipalities across the border
\item \textbf{Entry-restriction findings:} None; Uruguay's system is not a U.S.-style capped retail license
market
\item \textbf{Our contribution relative to this paper:} Different country, different outcome domain
(enforcement/seizures rather than legal retail market outcomes), and no entry-cap variation to test. Confirms
the general cross-border cannabis-policy spillover phenomenon extends outside the U.S., but does not bear on
our specific entry-regulation question.
\end{itemize}

\section{General industrial-organization entry literature (theoretical scaffolding, not cannabis-specific)}
Bresnahan and Reiss (1990, 1991), Berry (1992), Mazzeo (2002), Seim (2006), Toivanen and Waterson (2005), and
Berry and Reiss (2007) -- already cited and verified in our own manuscripts -- establish, in other retail and
service industries, that competitive conduct changes most sharply with the arrival of a market's second or
third entrant, and that the number of firms (not just their existence) disciplines price and behavior. None of
this literature has a cannabis application, and none pairs an entry-count mechanism with an exogenous demand
shock of the kind neighboring-state legalization timing provides. \textbf{Our contribution:} we apply this
general entry-threshold logic to a novel empirical setting (cannabis border markets) where the demand shock is
plausibly exogenous and precisely timed, rather than inferred from cross-sectional entrant-count variation
alone.

\section{Novelty assessment}
\textbf{No paper identified in this search directly tests whether cannabis entry restrictions moderate the
cross-border demand premium.} The entry-restriction literature (Thomas 2019; Hollenbeck and Uetake 2021) is
single-state (Washington) and has no demand-shock variation. The cross-border demand literature (Hansen,
Miller, and Weber 2020; Hao and Cowan 2020; Baggio, Chong, and Kwon 2020) treats the destination market's
regulatory structure as a fixed background fact, not a variable to test. Neither literature has been combined
with the other in the sources located. This is stated as a defensible, not an absolute, novelty claim -- see
Section~\ref{sec:methodology} for exactly what was and was not checked.

\section{Search methodology and coverage caveats}
\label{sec:methodology}
Approximately ten rounds of keyword searches were run using the term combinations specified for this review
(license cap, entry restriction, limited licenses, retail entry, border effect, border premium, cross-border
demand, border county, neighboring state, entry regulation combined with border), covering NBER, SSRN, and
RePEc/IDEAS listings, plus targeted checks of the \emph{RAND Journal of Economics}, \emph{Journal of Public
Economics}, \emph{Economic Inquiry}, and \emph{International Journal of Drug Policy}. One round of forward
citation-tracing was run on the two most central papers in each literature (Hansen, Miller, and Weber 2020 on
the border-demand side; Hollenbeck and Uetake 2021 and Thomas 2019 on the entry-restriction side).

\textbf{What was not done, stated plainly:} no direct access to Google Scholar or Web of Science; no
exhaustive state-by-state sweep beyond what surfaced organically (Arizona, Nevada, New Jersey, Maryland, and
Illinois were not individually checked for their own entry-cap literature); no verification of two
title-only results (Escudero; Kim, O'Connor, and Norwood) beyond their abstracts. Given this coverage, the
novelty conclusion above should be read as well-supported but not absolute -- exactly the standard the
original review request asked for.

\end{document}
